% ======================================================================
% Ontologie der Kontinua -- Core 1.3.3
% Cross-K-Modul: crossk_k1_k2.tex
% Cross-Level-Relationen zwischen K1 und K2
% ======================================================================

\subsubsection{K1–K2-Querverbindung}
\label{sec:crossk-k1-k2}

Der Übergang $K_1 \rightarrow K_2$ markiert das Auftreten geometrischer und
energetischer Struktur aus dem minimalen repräsentationalen Kontinuum.
Während $K_1$ eine einzelne Achse enthält, primitive Potentiale und triviale
Flüsse, führt $K_2$ ein:
\begin{itemize}
  \item mehrere Achsen, die eine Koordinatenstruktur bilden,
  \item explizite Potential--Gradienten-Beziehungen,
  \item nicht-triviale Flüsse $J_2$,
  \item Kausalordnung,
  \item und einen reicheren Domänenraum $\Omega(K_2)$, der physikalische
        Dynamik unterstützt.
\end{itemize}

Die Cross-Level-Struktur erfasst, wie die einfache repräsentationale Geometrie
von $K_1$ zu einer physikalisch ähnlichen Geometrie von $K_2$ erweitert wird.

% ----------------------------------------------------------------------
\subsubsection{1. Stufen und ihre Rollen}

\paragraph{K1.}
$K_1$ liefert:
\begin{itemize}
  \item eine einzelne Achse $A_1$,
  \item ein primitives Potential $P_1$,
  \item eine grundlegende Grenzstruktur $\partial\Omega(K_1)$,
  \item erste gradientenbasierte Flüsse $J_1$,
  \item triviale Zyklen, aber keine vollständige geometrische oder kausale Struktur.
\end{itemize}

\paragraph{K2.}
$K_2$ erweitert die Dimensionalität und führt ein:
\begin{itemize}
  \item mindestens eine neue Achse $A_2$,
  \item ein mehrdimensionales Koordinatenchart,
  \item energetische Potentiale und Felder,
  \item nicht-triviale dynamische Flüsse $J_2$,
  \item kausale Ordnung und physikalisch interpretierbare Nebenbedingungen,
  \item einen reicheren Zustandsraum $\Omega(K_2)$ mit nicht-trivialer Topologie.
\end{itemize}

$K_2$ ist das erste Kontinuum, das Perkolation, Feldausbreitung
und dynamisches Verhalten unterstützt.

% ----------------------------------------------------------------------
\subsubsection{2. Gemeinsame und vererbte Achsen sowie Schwellen}

\paragraph{Achsen.}
Die Achse $A_1$ von $K_1$ wird zu einer Koordinatenachse innerhalb von
$A(K_2)$:
\[
A(K_1) \subset A(K_2) = \{A_1, A_2, \dots\}.
\]
Die neue Achse $A_2$ ermöglicht das Auftreten von Gradienten, Feldern und
Transport über einen mehrdimensionalen Raum.

\paragraph{Schwellen.}
Die Existenzschwelle $\Theta_1$ garantiert die Stabilität von $A_1$; bei
Verletzung dieses Werts kollabiert der gesamte Übergang:
\[
\Theta_1^{\mathrm{death}} \Rightarrow \Omega(K_2)=\varnothing.
\]

$K_2$ führt neue Schwellenklassen ein:
\begin{itemize}
  \item \textbf{energetische Schwellen} für Felder und Potentiale,
  \item \textbf{Gradientenschwellen} für Flüsse,
  \item \textbf{Perkolationsschwellen} für Konnektivität,
  \item \textbf{Kausalitätsschwellen} für die Ordnung von Ereignissen.
\end{itemize}

Diese Schwellen verfeinern und erweitern die aus $K_1$ vererbte Struktur.

% ----------------------------------------------------------------------
\subsubsection{3. Cross-Level-Flüsse, Zyklen und Spannungen}

\paragraph{Flüsse.}
Flüsse $J_2$ gehen aus $J_1$ hervor, sind aber durch die zusätzliche Achse und
die energetischen Potentiale komplexer:
\[
J_2 = \nabla P_2 + \text{Transportterme entlang } A_2.
\]
Hierzu gehören diffusionsähnliches, feldartiges und kausales Ausbreitungsverhalten.

\paragraph{Zyklen.}
Zyklen in $K_2$ entsprechen geschlossenen geometrischen oder energetischen
Schleifen:
\[
C_2 : \oint_{\gamma} J_2 \cdot dA.
\]
Diese sind in $K_1$ unmöglich, da dort die Dimensionalität fehlt.
Somit ist $C_2$ eine strikte Aufwärts-Extension des trivialen Zyklus von $K_1$.

\paragraph{Spannung.}
Die Cross-Level-Spannung ist definiert als:
\[
T_{1,2} = h(P_2 - P_1,\ \Theta_2 - \Theta_1,\ A_2).
\]
Überschreitet $T_{1,2}$ die dimensionsbezogene Schwelle, wird die neue Achse
$A_2$ instabil und $\Omega(K_2)$ kollabiert zurück in eine degenerierte Form.

% ----------------------------------------------------------------------
\subsubsection{4. Geburts- und Todesbedingungen über die Stufen}

\paragraph{Geburt von \texorpdfstring{$K_2$}{K_2}.}
$K_2$ entsteht, wenn:
\[
T_1 > \Theta_{1,\mathrm{dim}}
\quad \text{und} \quad
A_2 \in M_2 \setminus A(K_1),
\]
mit entsprechender Erweiterung des Zustandsraums:
\[
\Omega(K_2) = \Omega(K_1) \cup \Delta\Omega_{\mathrm{geom}}.
\]

\paragraph{Todesausbreitung.}
Wenn $\Theta_1$ versagt, gibt es keine Aufwärts-Extension.
Wenn $\Theta_2$ versagt, kollabiert nur $K_2$; $K_1$ bleibt gültig, weil seine
Struktur nicht von den höherstufigen Potentialen oder Flüssen abhängt.

% ----------------------------------------------------------------------
\subsubsection{5. Konzeptionelle Beispiele}

\paragraph{Entstehung der Geometrie.}
Eine einzelne Achse ($K_1$) kann keine Flüsse mit Krümmung, Divergenz oder
Mehrfachausbreitung unterstützen. Die Ergänzung durch $A_2$ ermöglicht:
\[
\mathrm{grad},\ \mathrm{div},\ \mathrm{curl}-\text{artige Strukturen}.
\]

\paragraph{Perkolation.}
Ein mehrdimensionaler Raum unterstützt Konnektivitätsübergänge
(Perkolationsschwelle $p_c$), die Teil der $K_2$-Schwellenstruktur werden, aber
in $K_1$ kein Gegenstück haben.

\paragraph{Kausalität.}
Die gerichtete Ausbreitung über eine mehrdimensionale Struktur führt zum ersten
Auftreten von kausaler Ordnung: eine Unmöglichkeit in $K_1$.

Diese Beispiele veranschaulichen die Kernaussage des K1→K2-Übergangs:
\[
\textbf{Es ist die Geburt von Geometrie, Feldern und Dynamik aus dem minimalen
repräsentationalen Kontinuum.}
\]
